\documentclass{article}
\usepackage[spanish]{babel}
\usepackage[T1]{fontenc}
\usepackage[utf8]{inputenc}
\usepackage{amsmath}

\title{M\'etodos variacionales para ecuaciones diferenciales parciales}
\author{A.~Garc\'ia L\'opez}
\date{2026}

\begin{document}
\maketitle

\section{Introducci\'on}

Los m\'etodos variacionales constituyen una herramienta fundamental en el
an\'alisis de ecuaciones diferenciales parciales.  En este trabajo estudiamos
la existencia de soluciones d\'ebiles para problemas el\'ipticos no lineales
en dominios acotados de~$\mathbb{R}^n$.

% IB-001: Iberian vs Latin American mix
El ordenador realiza los c\'alculos num\'ericos con precisi\'on.  La computadora
procesa los datos de entrada.  El conductor del autobus facilit\'o la informaci\'on.
La informaci\'on fue transmitida por el celular.

% Correct: consistent Iberian Spanish
El ordenador ejecuta el algoritmo en tiempo polinomial.  El c\'alculo se realiza
mediante un programa inform\'atico desarrollado en la universidad.

% Inverted punctuation — correct usage
?`Cu\'al es la tasa de convergencia del m\'etodo propuesto?  !`Excelente resultado!
La respuesta es sorprendente.  ?`Podemos generalizar este enfoque a dimensiones
superiores?

% Missing inverted punctuation — incorrect
Cual es el error de aproximacion?  Es posible mejorar la cota!

\section{Formulaci\'on del problema}

Sea $\Omega \subset \mathbb{R}^n$ un dominio acotado con frontera regular.
Consideramos el problema
\begin{equation}
  -\Delta u + u^3 = f \quad \text{en } \Omega, \qquad u = 0 \quad \text{en } \partial\Omega.
\end{equation}

Definimos el funcional de energ\'ia $J \colon H_0^1(\Omega) \to \mathbb{R}$ como
\[
  J(u) = \frac{1}{2}\int_\Omega |\nabla u|^2\,dx + \frac{1}{4}\int_\Omega u^4\,dx
         - \int_\Omega fu\,dx.
\]

\section{Resultado principal}

\begin{theorem}
Si $f \in L^2(\Omega)$, entonces existe una \'unica soluci\'on d\'ebil
$u \in H_0^1(\Omega)$ del problema~(1).  Adem\'as, $u$ minimiza el funcional~$J$.
\end{theorem}

La demostraci\'on se basa en el teorema de Lax-Milgram y en la desigualdad
de Poincar\'e.  Utilizamos tambi\'en la coercividad del operador asociado.

\section{Experimentos num\'ericos}

Se implementaron elementos finitos de orden $p = 1, 2, 3$ en mallas triangulares.
Los resultados confirman la tasa de convergencia te\'orica de $O(h^{p+1})$ en la
norma~$L^2$.  La Tabla~1 resume los errores obtenidos.

\section{Conclusiones}

Hemos demostrado la existencia y unicidad de soluciones d\'ebiles para una clase
de problemas el\'ipticos no lineales.  El m\'etodo variacional proporciona un marco
natural para el an\'alisis y la aproximaci\'on num\'erica de estas ecuaciones.

\end{document}
